\documentclass[14pt]{extarticle} % Utilisation de extarticle pour supporter 14pt

% ------------------------------------------------------------------------
% Packages indispensables et recommandés
% ------------------------------------------------------------------------
\usepackage[utf8]{inputenc}     % Encodage des caractères
\usepackage[T1]{fontenc}        % Encodage de la police
\usepackage[french]{babel}      % Support de la langue française
\usepackage{amsmath, amssymb}   % Environnements mathématiques enrichis
\usepackage{amsthm}             % Environnements théorèmes
\usepackage{graphicx}           % Inclusion d'images
\usepackage{hyperref}           % Liens hypertextes
\usepackage{geometry}           % Gestion des marges
\usepackage{titlesec}           % Personnalisation des titres
\usepackage{setspace}           % Gestion de l'interligne
\usepackage{xcolor}             % Gestion des couleurs
\usepackage{booktabs}           % Pour améliorer les tableaux
\usepackage{fancyhdr}           % Personnalisation des en-têtes et pieds de page
\usepackage{cleveref}           % Références intelligentes
\usepackage{caption}            % Pour personnaliser les légendes
\usepackage{enumitem}           % Pour personnaliser les listes
\usepackage{pdflscape}          % Pour placer les tableaux en paysage

% ------------------------------------------------------------------------
% Configuration de la mise en page
% ------------------------------------------------------------------------
\geometry{
    a4paper,
    margin=25mm
}

% ------------------------------------------------------------------------
% Définition des environnements théorèmes, définitions, etc.
% ------------------------------------------------------------------------
\theoremstyle{definition}
\newtheorem{definition}{Définition}[section]

\theoremstyle{plain}
\newtheorem{proposition}{Proposition}[section]
\newtheorem{theorem}{Théorème}[section]

% ------------------------------------------------------------------------
% Personnalisation des sections
% ------------------------------------------------------------------------
\titleformat{\section}{\Large\bfseries}{\thesection}{1em}{}
\titleformat{\subsection}{\large\bfseries}{\thesubsection}{1em}{}
\titleformat{\subsubsection}{\normalsize\bfseries}{\thesubsubsection}{1em}{}

% ------------------------------------------------------------------------
% Commande pour mettre en avant les livres en bleu
% ------------------------------------------------------------------------
\newcommand{\livre}[1]{\textcolor{blue}{#1}}

% ------------------------------------------------------------------------
% Configuration des hyperliens
% ------------------------------------------------------------------------
\hypersetup{
    colorlinks=true,          % Les liens sont colorés
    linkcolor=blue,           % Couleur des liens internes (TOC, références, etc.)
    urlcolor=blue,            % Couleur des URL
    citecolor=blue,           % Couleur des citations
    filecolor=blue,           % Couleur des liens vers des fichiers
    pdfborder={0 0 0},        % Pas de bordure autour des liens
    breaklinks=true            % Permet les sauts de ligne dans les liens
}
% ------------------------------------------------------------------------
% Configuration des en-têtes et pieds de page
% ------------------------------------------------------------------------
\pagestyle{fancy}
\fancyhf{}
\fancyhead[L]{Transformation de Prais-Winsten}
\fancyhead[R]{Hiver 2025}
\fancyfoot[C]{\thepage}
\fancyfoot[R]{\includegraphics[height=1cm]{logo_universite_laval.png}} % Ajout du logo ERS en bas à droite


% ------------------------------------------------------------------------
% Informations sur le document
% ------------------------------------------------------------------------
\title{\textbf{Preuve Complète de la Transformation de Prais-Winsten}}
\author{GSF-6053}
\date{Janvier 2025}

% ------------------------------------------------------------------------
% Début du document
% ------------------------------------------------------------------------
\begin{document}
\maketitle

\tableofcontents

\newpage

\section{Introduction}
La transformation de Prais-Winsten est une méthode utilisée pour corriger l'autocorrélation des erreurs de premier ordre dans un modèle de régression linéaire. Cette transformation est particulièrement utile dans les séries temporelles où l'autocorrélation peut biaiser les estimations des moindres carrés ordinaires (OLS). La méthode de Prais-Winsten améliore la procédure de Cochrane-Orcutt en conservant la première observation, ce qui permet une utilisation plus efficace des données disponibles.

\section{Modèle avec Autocorrélation des Erreurs}
Considérons le modèle de régression linéaire suivant avec des erreurs autocorrélées d'ordre 1 :
\begin{align}
Y_t &= X_t \beta + u_t \label{eq:model_original} \\
u_t &= \rho u_{t-1} + \epsilon_t \label{eq:autocorr_error}
\end{align}
où \( \epsilon_t \) est un terme d'erreur aléatoire indépendant et identiquement distribué (\( \epsilon_t \sim \mathcal{N}(0, \sigma^2) \)) et \( |\rho| < 1 \).

\section{Objectif de la Transformation}
L'objectif principal de la transformation de Prais-Winsten est de transformer le modèle original de manière à éliminer l'autocorrélation des erreurs, permettant ainsi l'application des moindres carrés ordinaires pour obtenir des estimateurs sans biais et efficaces des paramètres \( \beta \).

\section{Transformation de Prais-Winsten}
La transformation de Prais-Winsten est définie par les équations suivantes pour \( t = 1, 2, \dots, T \) :
\begin{align}
Y_t^* &= Y_t - \rho Y_{t-1} \\
X_t^* &= X_t - \rho X_{t-1} \\
u_t^* &= u_t - \rho u_{t-1}
\end{align}
Contrairement à la méthode de Cochrane-Orcutt, Prais-Winsten modifie également la première observation pour éviter de perdre des données, ce qui rend l'estimation plus efficace.

\subsection{Procédure de Transformation}
\begin{enumerate}
    \item \textbf{Estimation Initiale :} Estimer le modèle OLS (\ref{eq:model_original}) pour obtenir les résidus \( \hat{u}_t \).
    \item \textbf{Estimation de \( \rho \) :} Utiliser les résidus pour estimer le coefficient d'autocorrélation \( \rho \) en régressant \( \hat{u}_t \) sur \( \hat{u}_{t-1} \).
    \begin{equation}
    \hat{u}_t = \rho \hat{u}_{t-1} + \nu_t
    \end{equation}
    \item \textbf{Transformation des Variables :} Appliquer la transformation de Prais-Winsten aux variables dépendante et indépendantes.
    \item \textbf{Estimation Finale :} Estimer le modèle transformé en utilisant les moindres carrés ordinaires sur les variables transformées.
\end{enumerate}

\section{Preuve de la Transformation}
Pour démontrer que la transformation de Prais-Winsten élimine l'autocorrélation des erreurs, nous allons procéder par étapes détaillées.

\subsection{Modèle Original avec Autocorrélation}
Considérons le modèle original avec autocorrélation des erreurs d'ordre 1 :
\begin{align}
Y_t &= X_t \beta + u_t \label{eq:model_original_repeat} \\
u_t &= \rho u_{t-1} + \epsilon_t \label{eq:autocorr_error_repeat}
\end{align}
où \( \epsilon_t \) est un terme d'erreur aléatoire indépendant avec \( E(\epsilon_t) = 0 \) et \( Var(\epsilon_t) = \sigma^2 \), et \( |\rho| < 1 \).

\subsection{Substitution et Transformation}
Substituons l'expression de \( u_t \) dans le modèle original (\ref{eq:model_original_repeat}) :
\begin{align*}
Y_t &= X_t \beta + \rho u_{t-1} + \epsilon_t \quad \text{(en remplaçant \( u_t \))} \\
Y_t &= X_t \beta + \rho (Y_{t-1} - X_{t-1} \beta) + \epsilon_t \quad \text{(en remplaçant \( u_{t-1} \))} \\
Y_t &= X_t \beta + \rho Y_{t-1} - \rho X_{t-1} \beta + \epsilon_t \quad \text{(en distribuant \( \rho \))}
\end{align*}
Regroupons les termes :
\begin{align*}
Y_t - \rho Y_{t-1} &= X_t \beta - \rho X_{t-1} \beta + \epsilon_t \\
Y_t^* &= X_t^* \beta + \epsilon_t \quad \text{(définissant \( Y_t^* = Y_t - \rho Y_{t-1} \) et \( X_t^* = X_t - \rho X_{t-1} \))}
\end{align*}
Ainsi, le modèle transformé est :
\[
Y_t^* = X_t^* \beta + \epsilon_t
\]
où \( Y_t^* = Y_t - \rho Y_{t-1} \) et \( X_t^* = X_t - \rho X_{t-1} \).

\subsection{Propriétés des Erreurs Transformées}
Examinons les erreurs transformées :
\[
u_t^* = u_t - \rho u_{t-1}
\]
En utilisant l'équation (\ref{eq:autocorr_error_repeat}) :
\[
u_t^* = (\rho u_{t-1} + \epsilon_t) - \rho u_{t-1} = \epsilon_t
\]
Ainsi, \( u_t^* = \epsilon_t \).

Puisque \( \epsilon_t \) est non corrélée et homoscédastique (\( Var(\epsilon_t) = \sigma^2 \)), les erreurs transformées \( u_t^* \) possèdent les mêmes propriétés :
\[
E(u_t^*) = E(\epsilon_t) = 0
\]
\[
Var(u_t^*) = Var(\epsilon_t) = \sigma^2
\]
\[
Cov(u_t^*, u_s^*) = Cov(\epsilon_t, \epsilon_s) = 0 \quad \text{pour } t \neq s
\]
Cela montre que les erreurs transformées \( u_t^* \) sont non corrélées et homoscédastiques.

\subsection{Conclusion de la Transformation}
En appliquant la transformation de Prais-Winsten, le modèle transformé devient :
\[
Y_t^* = X_t^* \beta + \epsilon_t
\]
où les erreurs \( \epsilon_t \) sont non corrélées et homoscédastiques. Cela permet l'application des moindres carrés ordinaires (OLS) sur le modèle transformé pour obtenir des estimateurs \( \hat{\beta} \) sans biais et efficaces.

\subsection{Étapes Détaillées de la Preuve}
Pour une compréhension approfondie, détaillons chaque étape de la preuve :

\subsubsection{Étape 1 : Substitution des Erreurs}
Partant des équations (\ref{eq:model_original_repeat}) et (\ref{eq:autocorr_error_repeat}) :
\[
Y_t = X_t \beta + u_t
\]
\[
u_t = \rho u_{t-1} + \epsilon_t
\]
Substituons \( u_t \) dans l'équation de \( Y_t \) :
\[
Y_t = X_t \beta + \rho u_{t-1} + \epsilon_t
\]

\subsubsection{Étape 2 : Substitution Récursive}
Remplaçons \( u_{t-1} \) par son expression à partir de l'équation (\ref{eq:autocorr_error_repeat}) :
\[
u_{t-1} = \rho u_{t-2} + \epsilon_{t-1}
\]
Cependant, pour simplifier la transformation, nous considérons que les valeurs passées de \( u \) sont déjà intégrées dans les termes de \( Y \) et \( X \).

\subsubsection{Étape 3 : Réarrangement des Termes}
Regroupons les termes liés à \( Y_t \) et \( Y_{t-1} \) :
\[
Y_t - \rho Y_{t-1} = X_t \beta - \rho X_{t-1} \beta + \epsilon_t
\]
\[
Y_t^* = X_t^* \beta + \epsilon_t
\]
avec \( Y_t^* = Y_t - \rho Y_{t-1} \) et \( X_t^* = X_t - \rho X_{t-1} \).

\subsubsection{Étape 4 : Analyse des Erreurs Transformées}
Les erreurs transformées \( u_t^* \) sont définies comme :
\[
u_t^* = u_t - \rho u_{t-1} = \epsilon_t
\]
Puisque \( \epsilon_t \) est indépendant et identiquement distribué, les erreurs transformées \( u_t^* \) sont également indépendantes et identiquement distribuées, sans autocorrélation.

\subsubsection{Étape 5 : Application des Moindres Carrés Ordinaires}
Avec le modèle transformé :
\[
Y_t^* = X_t^* \beta + \epsilon_t
\]
et des erreurs \( \epsilon_t \) non autocorrélées et homoscédastiques, l'application des moindres carrés ordinaires (OLS) sur ce modèle transformé fournit des estimateurs \( \hat{\beta} \) sans biais et efficaces, satisfaisant les conditions du théorème de Gauss-Markov.

\section{Estimation de \( \rho \)}
L'estimation du coefficient d'autocorrélation \( \rho \) est cruciale pour la transformation de Prais-Winsten. Une méthode courante consiste à utiliser les résidus du modèle OLS initial pour estimer \( \rho \) via une régression des résidus sur leurs valeurs retardées :
\[
\hat{u}_t = \rho \hat{u}_{t-1} + \nu_t
\]
où \( \nu_t \) est l'erreur de cette régression. L'estimateur de \( \rho \) obtenu est ensuite utilisé dans les étapes de transformation des variables.

\subsection{Procédure d'Estimation de \( \rho \)}
\begin{enumerate}
    \item \textbf{Estimation Initiale du Modèle OLS :} Estimer le modèle original (\ref{eq:model_original_repeat}) par OLS pour obtenir les résidus \( \hat{u}_t \).
    \item \textbf{Régression des Résidus :} Régressuer \( \hat{u}_t \) sur \( \hat{u}_{t-1} \) :
    \[
    \hat{u}_t = \rho \hat{u}_{t-1} + \nu_t
    \]
    \item \textbf{Obtention de \( \hat{\rho} \) :} Utiliser les estimateurs de cette régression pour obtenir \( \hat{\rho} \).
\end{enumerate}

\subsection{Considérations sur l'Estimation de \( \rho \)}
\begin{itemize}
    \item \textbf{Consistance :} L'estimateur \( \hat{\rho} \) est consistant si \( \rho \) est bien spécifié et si les erreurs sont correctement modélisées.
    \item \textbf{Efficacité :} Pour améliorer l'efficacité de l'estimation, des méthodes comme le maximum de vraisemblance peuvent être utilisées.
    \item \textbf{Tests de Significativité :} Il est important de tester la significativité de \( \hat{\rho} \) pour s'assurer que l'autocorrélation est effectivement présente et significative.
\end{itemize}

\section{Propriétés de l'Estimateur Transformé}
Après transformation, l'estimateur \( \hat{\beta} \) obtenu via les moindres carrés ordinaires sur le modèle transformé possède les propriétés suivantes :
\begin{itemize}
    \item \textbf{Sans Biais :} \( E(\hat{\beta}) = \beta \)
    \item \textbf{Efficace :} Sous les hypothèses classiques de Gauss-Markov, \( \hat{\beta} \) est l'estimateur linéaire sans biais ayant la plus petite variance.
    \item \textbf{Consistant :} L'estimateur converge en probabilité vers la vraie valeur \( \beta \) lorsque la taille de l'échantillon augmente.
\end{itemize}

\subsection{Démonstration de l'Unbiasedness}
Considérons le modèle transformé :
\[
Y_t^* = X_t^* \beta + \epsilon_t
\]
Prenons l'espérance conditionnelle de \( Y_t^* \) :
\[
E(Y_t^* | X_t^*) = E(X_t^* \beta + \epsilon_t | X_t^*) = X_t^* \beta + E(\epsilon_t | X_t^*) = X_t^* \beta
\]
Ainsi, l'estimateur \( \hat{\beta} \) issu des moindres carrés ordinaires est sans biais :
\[
E(\hat{\beta}) = \beta
\]

\subsection{Démonstration de l'Efficacité}
Sous les hypothèses de Gauss-Markov (linéarité, exogénéité, homoscédasticité, absence d'autocorrélation, etc.), l'estimateur OLS est BLUE (Best Linear Unbiased Estimator). Après transformation de Prais-Winsten, ces hypothèses sont satisfaites dans le modèle transformé, ce qui garantit que \( \hat{\beta} \) est l'estimateur linéaire sans biais avec la plus petite variance possible.

\subsection{Démonstration de la Consistance}
La consistance de \( \hat{\beta} \) repose sur les conditions suivantes :
\begin{itemize}
    \item \( \frac{1}{T} X^{*'} X^* \) converge vers une matrice définie positive.
    \item \( \frac{1}{T} X^{*'} \epsilon \) converge en probabilité vers zéro.
\end{itemize}
Sous ces conditions, par le théorème des moindres carrés, \( \hat{\beta} \) converge en probabilité vers \( \beta \).

\section{Conclusion}
La transformation de Prais-Winsten est une méthode efficace pour corriger l'autocorrélation des erreurs dans un modèle de régression linéaire. En transformant les variables dépendantes et indépendantes, cette méthode permet l'application des moindres carrés ordinaires de manière appropriée, assurant ainsi des estimations fiables des paramètres du modèle. Comparée à la méthode de Cochrane-Orcutt, Prais-Winsten conserve la première observation, ce qui améliore l'efficacité des estimations, surtout dans les petits échantillons.

\section{Références}
Pour une discussion détaillée sur la transformation de Prais-Winsten et ses propriétés, voir \textcolor{blue}{Wooldridge chapitre 8}.

\end{document}
